\documentclass[12pt,a4paper]{article}

% ==================== Packages ====================
\usepackage[
top=2.2cm,
bottom=2.3cm,
left=2.2cm,
right=2.2cm
]{geometry}

\usepackage{amsmath,amssymb}
\usepackage{graphicx}
\usepackage{float}
\usepackage{xcolor}
\usepackage{hyperref}
\usepackage{setspace}
\usepackage{enumitem}
\usepackage{titlesec}
\usepackage{caption}
\usepackage{tikz}
\usepackage[most]{tcolorbox}
\usepackage{xepersian}

% ==================== Fonts ====================
\settextfont{Yas}
\setdigitfont{Yas}
\setlatintextfont{Times New Roman}

% ==================== Colors ====================
\definecolor{ink}{RGB}{28,32,38}
\definecolor{navy}{RGB}{18,45,78}
\definecolor{teal}{RGB}{0,126,132}
\definecolor{muted}{RGB}{104,112,122}
\definecolor{paper}{RGB}{247,248,250}
\definecolor{line}{RGB}{220,225,231}

% ==================== General Style ====================
\onehalfspacing
\setlength{\parindent}{0pt}
\setlength{\parskip}{0.45em}

\hypersetup{
	colorlinks=true,
	linkcolor=teal,
	urlcolor=teal
}

\captionsetup{
	font=small,
	labelfont=bf,
	justification=centering
}

\setlist[itemize]{
	topsep=0.2em,
	itemsep=0.25em,
	rightmargin=1.4em
}

% ==================== Section Style ====================
\titleformat{\section}
{\Large\bfseries\color{navy}}
{\thesection}
{0.6em}
{}
[\vspace{-0.25em}{\color{teal}\titlerule[0.9pt]}]

\titleformat{\subsection}
{\large\bfseries\color{ink}}
{\thesubsection}
{0.6em}
{}

\titleformat{\subsubsection}
{\normalsize\bfseries\color{teal}}
{}
{0pt}
{}

% ==================== Boxes ====================
\tcbset{
	sharp corners,
	boxrule=0pt,
	left=9pt,
	right=9pt,
	top=8pt,
	bottom=8pt,
	colback=paper,
	colframe=paper
}

\newtcolorbox{summarybox}{
	colback=paper,
	borderline west={3pt}{0pt}{teal}
}

\newtcolorbox{resultbox}{
	colback=white,
	borderline west={3pt}{0pt}{navy},
	boxrule=0.5pt,
	colframe=line
}

% ==================== Document ====================
\begin{document}
	
	% ==================== Cover ====================
	\begin{titlepage}

		
		\vspace*{1.7cm}
		
		{\Huge\bfseries\color{navy} گزارش پروژه سوم\par}
		\vspace{0.35cm}
		{\Large\bfseries نرم‌افزار ریاضی\par}
		
		\vspace{1.1cm}
		
		{\color{teal}\rule{0.58\textwidth}{1.4pt}}
		
		\vspace{1.2cm}
		
		\begin{summarybox}
			\Large
			این گزارش شامل پیاده‌سازی و تحلیل سه مسئله عددی است:  
			درونیابی، تحلیل زمان اجرای الگوریتم‌ها و محاسبه عددی انتگرال.
		\end{summarybox}
		
		\vspace{1.4cm}
		
		\begin{tabular}{rl}
			\textbf{نام دانشجو:} & محمد جواد مسگرها \\[0.35cm]
			\textbf{شماره دانشجویی:} & 40312043 \\[0.35cm]
			\textbf{موضوع:} & \lr{Numerical Methods} \\[0.35cm]
			\textbf{مخزن پروژه:} &
\href{https://github.com/JaVaD15Mesgarha/mathematics_softwares_project}
{\lr{\texttt{GitHub Repository}}}

		\end{tabular}
		
		\vfill
		
		{\color{muted}\today}
		
	\end{titlepage}
	
	\tableofcontents
	\newpage
	
	% ==================== Intro ====================
	\section*{چکیده گزارش}
	\addcontentsline{toc}{section}{چکیده گزارش}
	
	در این پروژه، سه بخش اصلی از محاسبات عددی بررسی شده است. ابتدا دو روش درونیابی لاگرانژ و اسپلاین مکعبی روی یک مجموعه داده مقایسه می‌شوند. سپس زمان اجرای سه الگوریتم با استفاده از داده‌های فایل اکسل تحلیل می‌شود. در پایان نیز انتگرال تابع
	$e^{x^2}$
	با روش‌های ذوزنقه‌ای، سیمپسون و گاوس-لژاندر محاسبه و از نظر دقت مقایسه می‌شود.
	
	% ==================== Question 1 ====================
	\section{درونیابی لاگرانژ و اسپلاین مکعبی}
	
	\begin{summarybox}
		هدف این بخش، ساخت تابعی است که از نقاط داده‌شده عبور کند و رفتار دو روش لاگرانژ و اسپلاین مکعبی از نظر نوسان و پایداری مقایسه شود.
	\end{summarybox}
	
	نقاط مورد بررسی عبارت‌اند از:
	\[
	(1,1),\ (2,3),\ (3,5),\ (4,8),\ (5,5),\ (6,2)
	\]
	
	برای این داده‌ها، روش لاگرانژ یک چندجمله‌ای سراسری می‌سازد، اما اسپلاین مکعبی تابع را به‌صورت تکه‌ای و با چندجمله‌ای‌های درجه سه مدل می‌کند. هر دو روش در گره‌های اصلی خطای صفر دارند.
	
	\subsection{خروجی عددی}
	
	چندجمله‌ای لاگرانژ حاصل به صورت زیر است:
	\begin{equation}
		P(x)=
		0.175x^5
		-2.958x^4
		+18.375x^3
		-52.041x^2
		+68.45x
		-31.0
	\end{equation}
	
	در اسپلاین مکعبی، برای هر بازه یک تابع جداگانه ساخته می‌شود. نمونه قطعه اول در بازه
	$[1,2]$
	برابر است با:
	\begin{equation}
		S_0(x)=
		-0.186x^3
		+0.559x^2
		+1.626x
		-1.0
	\end{equation}
	
	\subsection{تحلیل}
	
	\begin{resultbox}
		لاگرانژ به دلیل سراسری بودن، در برخی نواحی مخصوصاً نزدیک انتهای بازه نوسان بیشتری دارد. در مقابل، اسپلاین مکعبی به دلیل ساختار تکه‌ای، منحنی نرم‌تر و کنترل‌شده‌تری تولید می‌کند. بنابراین برای داده‌های تجربی، اسپلاین انتخاب مطمئن‌تری است.
	\end{resultbox}
	
	\begin{figure}[H]
		\centering
		\includegraphics[width=0.84\textwidth]{../outputs/figures/question1_interpolation_v2.png}
		\caption{مقایسه درونیابی لاگرانژ و اسپلاین مکعبی}
	\end{figure}
	
	\newpage
	
	% ==================== Question 2 ====================
	\section{تحلیل زمان اجرای الگوریتم‌ها}
	
	\begin{summarybox}
		در این بخش زمان اجرای سه الگوریتم برای داده‌هایی با اندازه
		$100$
		تا
		$700\,\text{KB}$
		بررسی شده و رفتار آن‌ها با نمودارهای خطی، میله‌ای و جعبه‌ای تحلیل می‌شود.
	\end{summarybox}
	
	داده‌ها از فایل اکسل خوانده شده‌اند و پیاده‌سازی به‌گونه‌ای انجام شده که با اضافه شدن داده‌های جدید، محاسبات و نمودارها بدون تغییر اساسی به‌روزرسانی شوند.
	
	\subsection{روش کار}
	
	داده‌ها با
	\lr{pandas.read\_excel}
	وارد برنامه شدند. سپس برای هر الگوریتم، شاخص‌های آماری شامل میانگین، میانه و انحراف معیار محاسبه شد. برای نمایش بهتر، سه نمودار تولید شد:
	
	\begin{itemize}
		\item نمودار خطی برای بررسی روند رشد زمان اجرا
		\item نمودار میله‌ای برای مقایسه مستقیم الگوریتم‌ها
		\item نمودار جعبه‌ای برای تحلیل پراکندگی داده‌ها
	\end{itemize}
	
	\subsection{مقایسه الگوریتم‌ها}
	
	\begin{resultbox}
		\textbf{\lr{Alg.1}} سریع‌ترین و پایدارترین الگوریتم است. زمان اجرای آن از
		$50\,\text{ms}$
		تا
		$80\,\text{ms}$
		تغییر می‌کند و رشد بسیار آرامی دارد.
	\end{resultbox}
	
	\begin{resultbox}
		\textbf{\lr{Alg.2}} در ابتدا کندتر است، اما شیب رشد آن کنترل‌شده باقی می‌ماند. این الگوریتم برای ورودی‌های متوسط و بزرگ رفتار قابل قبولی دارد.
	\end{resultbox}
	
	\begin{resultbox}
		\textbf{\lr{Alg.3}} در ورودی‌های کوچک عملکرد نسبتاً مناسبی دارد، اما با افزایش حجم داده، زمان اجرای آن سریع رشد می‌کند و در ورودی‌های بزرگ ضعیف‌تر ظاهر می‌شود.
	\end{resultbox}
	
	\subsection{نمودارها}
	
	نمودار خطی نشان می‌دهد که
	\lr{Alg.1}
	در تمام اندازه‌ها بهترین عملکرد را دارد. همچنین تقاطع
	\lr{Alg.2}
	و
	\lr{Alg.3}
	نشان می‌دهد که انتخاب الگوریتم می‌تواند به اندازه ورودی وابسته باشد.
	
	\begin{figure}[H]
		\centering
		\includegraphics[width=0.84\textwidth]{../outputs/figures/question2_line_chart.png}
		\caption{نمودار خطی زمان اجرای الگوریتم‌ها}
	\end{figure}
	
	\begin{figure}[H]
		\centering
		\includegraphics[width=0.84\textwidth]{../outputs/figures/question2_bar_chart.png}
		\caption{نمودار میله‌ای مقایسه زمان اجرا}
	\end{figure}
	
	\begin{figure}[H]
		\centering
		\includegraphics[width=0.84\textwidth]{../outputs/figures/question2_box_plot.png}
		\caption{نمودار جعبه‌ای پراکندگی زمان اجرا}
	\end{figure}
	
	\subsection{میانگین زمان اجرای الگوریتم دوم}
	
	برای بازه
	$100$
	تا
	$600\,\text{KB}$
	مقادیر زمان اجرای
	\lr{Alg.2}
	برابرند با:
	\[
	200,\ 220,\ 240,\ 260,\ 280,\ 300
	\]
	
	بنابراین:
	\[
	\text{Average}
	=
	\frac{200+220+240+260+280+300}{6}
	=
	250.0\,\text{ms}
	\]
	
	\newpage
	
	% ==================== Question 3 ====================
	\section{محاسبه عددی انتگرال}
	
	\begin{summarybox}
		در این بخش مقدار انتگرال
		$\int_0^1 e^{x^2}\,dx$
		با سه روش عددی محاسبه شده و خطای هر روش نسبت به مقدار مرجع بررسی می‌شود.
	\end{summarybox}
	
	مقدار مرجع انتگرال برابر است با:
	\[
	1.462651745907
	\]
	
	این مقدار با استفاده از
	\lr{scipy.integrate.quad}
	به دست آمده و مبنای محاسبه خطای مطلق قرار گرفته است.
	
	\subsection{روش‌های استفاده‌شده}
	
	سه روش زیر پیاده‌سازی شدند:
	
	\begin{itemize}
		\item روش ذوزنقه‌ای مرکب
		\item روش سیمپسون یک‌سوم مرکب
		\item روش گاوس-لژاندر
	\end{itemize}
	
	\subsection{تحلیل دقت}
	
	\begin{resultbox}
		روش ذوزنقه‌ای ساده‌ترین روش است، اما مرتبه خطای آن
		$O(h^2)$
		است و برای رسیدن به دقت بالا به تعداد زیربازه بیشتری نیاز دارد.
	\end{resultbox}
	
	\begin{resultbox}
		روش سیمپسون با مرتبه خطای
		$O(h^4)$
		دقیق‌تر از ذوزنقه‌ای است و با تعداد تقسیم‌بندی کمتر، خطای بسیار پایین‌تری تولید می‌کند.
	\end{resultbox}
	
	\begin{resultbox}
		روش گاوس-لژاندر بهترین عملکرد را دارد. چون نقاط و وزن‌ها بهینه انتخاب می‌شوند، برای تابعی هموار مانند
		$e^{x^2}$
		با تعداد کمی ارزیابی تابع به دقت بسیار بالا می‌رسد.
	\end{resultbox}
	
	\subsection{نتیجه}
	
	برای این مسئله، روش گاوس-لژاندر دقیق‌ترین و کارآمدترین روش است. اگر داده‌ها به‌صورت جدولی و با فاصله‌های مساوی باشند، سیمپسون گزینه مناسبی است. روش ذوزنقه‌ای نیز از نظر آموزشی و پیاده‌سازی ساده است، اما از نظر دقت ضعیف‌تر عمل می‌کند.
	
	\begin{figure}[H]
		\centering
		\includegraphics[width=0.84\textwidth]{../outputs/figures/question3_integration.png}
		\caption{نمایش تابع و مقایسه خطای روش‌های انتگرال‌گیری}
	\end{figure}
	
\end{document}
